\section{Architecture Decisions}
\label{sec:ArchitectureDecisions}

Durga is not only a generator. It is a generator with a fairly opinionated runtime shape.
Several choices were made deliberately because they improve inspectability and make the
system easier to discuss, test and evolve.

\subsection{Canonical lifecycle stream}

The most important decision is the shared \texttt{process-events} topic. Execution topics
carry work between generated handlers, but they are not treated as the authoritative query
surface. Durga instead emits normalized lifecycle events and builds state projections from
that stream.

This was chosen because:
\begin{itemize}
\item execution topics represent transport history, not current truth,
\item process instance state must be queryable independently of routing topology,
\item counts, latency summaries and stuck-instance views are easier to derive from explicit
lifecycle events than from task-topic inspection.
\end{itemize}

\subsection{Visible generated runtime}

Durga generates many explicit classes instead of one hidden execution engine. That is a
tradeoff: the output is more verbose, but it is also far easier to inspect. This supports
the project goal of making BPMN-to-Kafka mapping concrete rather than magical.

\subsection{Generated helper scripts as product surface}

The helper scripts are deliberate, not incidental. A BPMN scaffolding tool is difficult to
evaluate if the user must write custom producers and consumers before seeing anything happen.
Durga therefore generates probes, scenario runners and manual interaction scripts together
with the runtime.

\subsection{Kafka Streams as the default read model}

The monitoring side uses Kafka Streams because it fits the core question well: consume one
canonical stream and materialize queryable state from it. The current implementation is not
trying to solve every production reporting problem; it is solving the local and architectural
read-model problem directly.

\subsection{Scope-aware subprocess handling}

Subprocess support could have been left as simple flattening. Durga started that way, but
then moved toward explicit subprocess entry and completion services because scope behavior
matters once cancellation, failure, escalation and event subprocesses are introduced.
